% 図：ジョブの状態の切り替え（chapters/commandline/process.tex）
\begin{tikzpicture}[
    >=Stealth, auto, node distance=26mm and 50mm, on grid,
    every state/.style={draw, rounded corners, rectangle, fill=blue!6,
      minimum width=28mm, minimum height=10mm, font=\small, align=center}]
  \node[state, fill=blue!15] (fg) {\textbf{フォアグラウンド}\\{\footnotesize 端末を占有して実行}};
  \node[state, right=of fg, fill=green!8] (bg) {\textbf{バックグラウンド}\\{\footnotesize 裏で実行}};
  \node[state, below right=26mm and 25mm of fg, fill=orange!10] (st) {\textbf{停止中}\\{\footnotesize 一時停止}};
  \path[->, thick]
    (bg) edge node[font=\footnotesize, swap]{\texttt{fg}} (fg)
    (fg) edge node[font=\footnotesize]{\texttt{Ctrl+Z}} (st)
    (st) edge[bend right=20] node[font=\footnotesize, swap]{\texttt{bg}} (bg)
    (st) edge[bend left=40] node[font=\footnotesize]{\texttt{fg}} (fg);
  \node[font=\footnotesize, anchor=south] (s1) at ([yshift=6mm]fg.north) {\texttt{コマンド}};
  \node[font=\footnotesize, anchor=south] (s2) at ([yshift=6mm]bg.north) {\texttt{コマンド \&}};
  \draw[->, thick] (s1) -- (fg);
  \draw[->, thick] (s2) -- (bg);
\end{tikzpicture}
